\chapter*{Notation}
\addcontentsline{toc}{chapter}{Notation}
\label{front:notation}

\begin{center}
\begin{tabular}{@{}ll@{}}
\(\C,\R,\Z,\Q,\F_p\) & complex, real, integer, rational, and finite fields\\
\(\HH\) & upper half-plane\\
\(\Lambda\) & lattice \(\Z\omega_1+\Z\omega_2\)\\
\(T=\C/\Lambda\) & complex torus\\
\(O\) & identity of \(T\), also point at infinity of \(E\)\\
\(\wp,\wp'\) & Weierstrass function and derivative\\
\(g_2,g_3,\Delta\) & invariants and discriminant \(g_2^3-27g_3^2\)\\
\(E\) & cubic \(Y^2Z=4X^3-g_2XZ^2-g_3Z^3\)\\
\(\Div(X)\) & divisor group of \(X\)\\
\(\divisor(f)\) & principal divisor of \(f\)\\
\(\Pic^0(X)\) & degree-zero divisor classes modulo principal divisors\\
\(\smap\) & Abel--Jacobi sum map\\
\(L(D),\ell(D)\) & meromorphic functions bounded by \(D\), and its dimension\\
\(\Phi\) & Weierstrass map \([\wp:\wp':1]\)\\
\(\oplus\) & transported group law on \(E\)\\
\(\boxplus\) & chord--tangent construction\\
\(j\) & \(j\)-invariant \(1728g_2^3/\Delta\)\\
\end{tabular}
\end{center}

\noindent
The symbols in this table are used throughout the book.  Later chapters refine
some definitions, but this front-matter list fixes the basic notation before the
main exposition begins.
